% !TEX program = pdflatex
%==============================================================================
%  Rapport technique : reconstruction pédagogique et évaluation critique
%  du modèle cosmologique Janus de Petit, Margnat & Zejli (EPJC 2024).
%
%  Compilation : pdflatex (2 passes pour la table des matières et les renvois).
%  Sous Overleaf : régler le compilateur sur pdfLaTeX.
%==============================================================================
\documentclass[11pt,a4paper]{article}

%------------------------------ Encodage / langue ----------------------------
\usepackage[utf8]{inputenc}
\usepackage[T1]{fontenc}
\usepackage[french]{babel}
\usepackage{lmodern}

%------------------------------ Mathématiques --------------------------------
\usepackage{amsmath,amssymb,amsthm,mathtools}
\usepackage{bm}
\usepackage{physics}        % \dd, \dv, \pdv, \abs, \norm ... (présent sur Overleaf)

%------------------------------ Mise en page ---------------------------------
\usepackage[a4paper,margin=2.4cm]{geometry}
\usepackage{enumitem}
\usepackage{microtype}
\usepackage{booktabs}
\usepackage{array}
\usepackage{xcolor}
\usepackage[colorlinks=true,linkcolor=blue!55!black,
            citecolor=blue!55!black,urlcolor=teal!60!black]{hyperref}

%------------------------------ Théorèmes ------------------------------------
\theoremstyle{definition}
\newtheorem{definition}{Définition}[section]
\newtheorem{lemme}{Lemme}[section]
\newtheorem{remarque}{Remarque}[section]
\theoremstyle{plain}
\newtheorem{theoreme}{Théorème}[section]
\newtheorem{propriete}{Propriété}[section]

%------------------------------ Raccourcis -----------------------------------
\newcommand{\dx}{\,\dd^4x}
\newcommand{\gbar}{\bar{g}}
\newcommand{\Rbar}{\bar{R}}
\newcommand{\rhobar}{\bar{\rho}}
\newcommand{\abar}{\bar{a}}
\newcommand{\cbar}{\bar{c}}
\newcommand{\nablabar}{\bar{\nabla}}
\newcommand{\Tp}{T'}
\newcommand{\Tbar}{\bar{T}}
\newcommand{\Sbar}{\bar{S}}
\newcommand{\Tbarp}{\bar{T}'}
\newcommand{\chibar}{\bar{\chi}}
\newcommand{\sgn}{\operatorname{sgn}}
\newcommand{\diag}{\operatorname{diag}}
\renewcommand{\Tr}{\operatorname{Tr}}
\newcommand{\Lor}{\mathrm{Lor}}
\newcommand{\Poi}{\mathrm{Poi}}
\newcommand{\half}{\tfrac{1}{2}}
\newcommand{\eqpaper}[1]{\textsf{(éq.~#1)}}   % renvoi aux équations du papier
\DeclareMathOperator{\Ad}{Ad}

\setlength{\parskip}{0.4em}
\setlength{\parindent}{0pt}

%==============================================================================
\title{\bfseries Reconstruction pédagogique et évaluation critique\\
du modèle cosmologique Janus\\[4pt]
\large Petit, Margnat \& Zejli, \emph{Eur. Phys. J. C} \textbf{84}:1226 (2024)\\
\normalsize DOI~\href{https://doi.org/10.1140/epjc/s10052-024-13569-w}{10.1140/epjc/s10052-024-13569-w}}
\author{Rapport technique}
\date{\today}

%==============================================================================
\begin{document}
\maketitle

\begin{abstract}
\noindent
Ce rapport reconstruit \emph{ex nihilo} les outils mathématiques et physiques nécessaires
pour comprendre et redémontrer les calculs de l'article de Petit, Margnat \& Zejli
(ci-après \textbf{PMZ24}). Il s'adresse à un lecteur ayant un fort bagage en mathématiques
et en mécanique (EDP, calcul variationnel, algèbre linéaire) mais sans formation préalable
en relativité ni en géométrie différentielle. La progression part du calcul tensoriel et
de la relativité restreinte (Partie~\ref{part:fond}), introduit la géométrie symplectique
de Souriau et les groupes dynamiques (Partie~\ref{part:souriau}), résume l'argument
topologique (Partie~\ref{part:topo}), reconstruit en détail les équations de champ
bimétriques et leurs solutions FLRW et de Schwarzschild (Partie~\ref{part:champ}), puis
expose une évaluation critique sourcée (Partie~\ref{part:critique}), incluant l'objection
centrale de T.~Damour. Les renvois \eqpaper{X} pointent vers les équations de l'article
publié. Le statut éditorial est arrêté à mai~2026.

\medskip
\noindent\textbf{Avertissement de neutralité.} Une part de ce rapport présente des outils
de physique \emph{établie} (relativité restreinte et générale, géométrie symplectique) ;
une autre part présente des constructions \emph{spécifiques} au modèle Janus, qui sont
scientifiquement \emph{contestées}. Le texte distingue systématiquement les deux.
\end{abstract}

\tableofcontents
\bigskip\hrule\bigskip

%==============================================================================
\section*{Conventions et notations}
\addcontentsline{toc}{section}{Conventions et notations}

\textbf{Signature.} On adopte la signature $(+,-,-,-)$ du papier, soit la matrice de Gram
\begin{equation}
G=\eta=\diag(1,-1,-1,-1).
\label{eq:gram}
\end{equation}
C'est la convention « Souriau / physique des particules », opposée à la convention « MTW »
$(-,+,+,+)$. Le carré pseudo-euclidien d'un quadrivecteur $P$ vaut
$P^{\mathsf T}\eta\,P=E^2/c^2-\abs{\vb p}^2=m^2c^2$.

\textbf{Indices.} Grecs $\mu,\nu,\rho,\sigma=0,1,2,3$ (espace-temps) ; latins $i,j,k=1,2,3$
(spatiaux). Convention de sommation d'Einstein implicite. Barres ($\gbar,\Rbar,\nablabar,\dots$)
pour le second feuillet.

\textbf{Constante d'Einstein.} Deux conventions coexistent (cf.~\S\ref{sec:einstein}) :
$\chi_{\text{MTW}}=+8\pi G/c^4$ et $\chi_{\text{Souriau}}=-8\pi G/c^2$
(valeur numérique $1.856\times10^{-27}\,\mathrm{cm\,g^{-1}}$, \eqpaper{57}).

%==============================================================================
\part*{}
\section{Partie I — Fondations mathématiques et physiques}
\label{part:fond}

\subsection{Calcul tensoriel et géométrie différentielle élémentaire}

\begin{definition}[Vecteurs et covecteurs]
Sur un espace vectoriel $V$ de dimension $n$, un \emph{vecteur} $v\in V$ a des composantes
$v^a$ qui se transforment par une matrice $A$ : $v'^a=A^a{}_b\,v^b$. Un \emph{covecteur}
(forme linéaire) $\omega\in V^\ast$ a des composantes $\omega_a$ qui se transforment par
l'inverse transposée : $\omega'_a=(A^{-1})^b{}_a\,\omega_b$.
\end{definition}

D'où la convention de Ricci : indice \emph{haut} = contravariant (vecteur),
indice \emph{bas} = covariant (covecteur). Un tenseur de type $(p,q)$ est une application
multilinéaire $T:(V^\ast)^p\times V^q\to\mathbb R$, de loi de transformation
\begin{equation}
T'^{a_1\cdots a_p}{}_{b_1\cdots b_q}
= A^{a_1}{}_{c_1}\!\cdots A^{a_p}{}_{c_p}\,
(A^{-1})^{d_1}{}_{b_1}\!\cdots (A^{-1})^{d_q}{}_{b_q}\,
T^{c_1\cdots c_p}{}_{d_1\cdots d_q}.
\end{equation}

\begin{definition}[Métrique]
Une métrique $g$ est un tenseur $(0,2)$ symétrique non dégénéré. Elle induit un isomorphisme
$V\!\to\!V^\ast$ permettant de monter/descendre les indices :
$v_\mu=g_{\mu\nu}v^\nu$, $\omega^\mu=g^{\mu\nu}\omega_\nu$, où $g^{\mu\nu}$ est l'inverse
matriciel ($g^{\mu\rho}g_{\rho\nu}=\delta^\mu_\nu$). L'élément de longueur est
$\dd s^2=g_{\mu\nu}\dd x^\mu\dd x^\nu$.
\end{definition}

En espace plat de Minkowski : $\dd s^2=\dd t^2-\dd x^2-\dd y^2-\dd z^2$ \eqpaper{6}.

\subsection{Relativité restreinte et groupe de Lorentz}

\textbf{Postulat.} L'invariance de $\dd s^2$ sous un changement de référentiel inertiel
$x'^\mu=\Lambda^\mu{}_\nu x^\nu+C^\mu$ impose
\begin{equation}
L^{\mathsf T}\,G\,L=G \qquad\eqpaper{3},
\label{eq:lorentz}
\end{equation}
qui définit \emph{axiomatiquement} le groupe de Lorentz $\Lor$ (groupe de Lie de dimension~6 :
3~rotations + 3~boosts).

\begin{propriete}[Quatre composantes connexes]
$\Lor$ n'est pas connexe. Les invariants discrets $\det L=\pm1$ et $\sgn(L^0{}_0)$ découpent
quatre composantes \eqpaper{7} :
\[
\begin{array}{llll}
\Lor_n: & \det=+1,\ L^0{}_0\ge1 & \text{(neutre, contient } \mathbb 1) & \omega(+1,+1)\\
\Lor_s: & \det=-1,\ L^0{}_0\ge1 & \text{(parité } P) & \omega(+1,-1)\\
\Lor_t: & \det=-1,\ L^0{}_0\le-1 & \text{(renversement } T) & \omega(-1,+1)\\
\Lor_{st}: & \det=+1,\ L^0{}_0\le-1 & (PT) & \omega(-1,-1)
\end{array}
\]
avec le paramétrage compact $\omega(\lambda_1,\lambda_2)=\diag(\lambda_1,\lambda_2,\lambda_2,\lambda_2)$,
$\lambda_i\in\{-1,+1\}$, et $L=\omega(\lambda_1,\lambda_2)\,L_n$ \eqpaper{15--16}.
\end{propriete}

Sous-groupes \eqpaper{8--9} : orthochrone $\Lor_o=\Lor_n\cup\Lor_s$ (préserve le temps),
antichrone $\Lor_a=\Lor_t\cup\Lor_{st}$ (renverse le temps).

\textbf{Groupe de Poincaré} $\Poi=\Lor\ltimes\mathbb R^{1,3}$, représenté \eqpaper{2} par
\begin{equation}
g=\begin{pmatrix} L & C\\ 0 & 1\end{pmatrix},\qquad
g\cdot X = L\,X + C \quad\eqpaper{12--13},
\end{equation}
groupe de dimension~10.

\begin{theoreme}[Définition de la masse selon Souriau]
\label{th:masse}
Pour une particule massive, $P=(E/c,\vb p)$ est un quadrivecteur, $P'=LP$ \eqpaper{25}, et
$P^{\mathsf T}GP=E^2/c^2-\abs{\vb p}^2=m^2c^2$. Souriau \emph{définit} la masse avec le signe
de l'énergie \eqpaper{30} :
\begin{equation}
\boxed{\,m=\sqrt{P^{\mathsf T}G\,P}\;\sgn(E)\,}.
\label{eq:masse}
\end{equation}
\end{theoreme}

La conséquence-clé : si une transformation renverse $E\to-E$, alors $m\to-m$. C'est la
\emph{justification mathématique} (établie dans le formalisme symplectique) des masses
négatives ; son \emph{usage cosmologique} est en revanche spécifique à Janus et contesté.

\textbf{Symétries discrètes.} $P$ (parité, $\lambda_2=-1$) sans difficulté (deux hélicités du
photon). $T$ (renversement temporel, $\lambda_1=-1$) : pivot de Janus, inverse $E$.
$C$ (conjugaison de charge) : interprétée géométriquement via une $5^{\text e}$ dimension
(Partie~\ref{part:souriau}). \textbf{CPT} : théorème de Lüders–Pauli (toute théorie locale,
Lorentz-invariante, hermitienne est $CPT$-invariante). Sakharov en tire ses trois conditions
de baryogenèse \cite{Sakharov1967}.

\subsection{Géométrie différentielle pseudo-riemannienne}

\begin{definition}[Connexion de Levi-Civita]
L'unique connexion sans torsion compatible avec $g$ ($\nabla_\rho g_{\mu\nu}\equiv0$) a pour
coefficients les symboles de Christoffel
\begin{equation}
\Gamma^\rho{}_{\mu\nu}=\half\,g^{\rho\sigma}
\bigl(\partial_\mu g_{\sigma\nu}+\partial_\nu g_{\sigma\mu}-\partial_\sigma g_{\mu\nu}\bigr).
\label{eq:christoffel}
\end{equation}
Ce ne sont pas des composantes de tenseur.
\end{definition}

\textbf{Dérivée covariante} : $\nabla_\mu V^\nu=\partial_\mu V^\nu+\Gamma^\nu{}_{\mu\rho}V^\rho$,
$\nabla_\mu\omega_\nu=\partial_\mu\omega_\nu-\Gamma^\rho{}_{\mu\nu}\omega_\rho$.

\textbf{Géodésique} : $\ddot x^\mu+\Gamma^\mu{}_{\nu\rho}\dot x^\nu\dot x^\rho=0$.

\textbf{Courbure.} Le tenseur de Riemann mesure le défaut de commutation :
$[\nabla_\mu,\nabla_\nu]V^\rho=R^\rho{}_{\sigma\mu\nu}V^\sigma$, avec
\begin{equation}
R^\rho{}_{\sigma\mu\nu}=\partial_\mu\Gamma^\rho{}_{\nu\sigma}-\partial_\nu\Gamma^\rho{}_{\mu\sigma}
+\Gamma^\rho{}_{\mu\lambda}\Gamma^\lambda{}_{\nu\sigma}
-\Gamma^\rho{}_{\nu\lambda}\Gamma^\lambda{}_{\mu\sigma}.
\end{equation}
Ricci : $R_{\mu\nu}=R^\rho{}_{\mu\rho\nu}$ ; scalaire : $R=g^{\mu\nu}R_{\mu\nu}$.

\begin{theoreme}[Identités de Bianchi contractées]
\label{th:bianchi}
\begin{equation}
\nabla_\mu\!\left(R^{\mu\nu}-\half R\,g^{\mu\nu}\right)=0.
\label{eq:bianchi}
\end{equation}
Cette identité garantit la conservation du membre de gauche d'Einstein et impose
$\nabla_\mu T^{\mu\nu}=0$.
\end{theoreme}

\begin{remarque}[Rôle dans la critique de Damour]
Avec \emph{deux} métriques $(g,\gbar)$ on a \emph{deux} connexions $(\nabla,\nablabar)$ et donc
\emph{deux} identités de Bianchi indépendantes. Si la \emph{même} matière source les deux
équations, elle est soumise à deux lois de conservation potentiellement incompatibles.
C'est le cœur de l'objection de \S\ref{sec:damour}.
\end{remarque}

\subsection{Équations d'Einstein et dérivation variationnelle}
\label{sec:einstein}

\textbf{Équations d'Einstein (1915)} :
\begin{equation}
R_{\mu\nu}-\half R\,g_{\mu\nu}=\chi\,T_{\mu\nu}.
\label{eq:einstein}
\end{equation}

\textbf{Conventions de $\chi$.}
\begin{itemize}[leftmargin=1.4em,nosep]
\item MTW / Wald / Damour : $\chi=+8\pi G/c^4$, $T^{00}=\rho c^2$.
\item Souriau / Petit \eqpaper{57} : $\chi=-8\pi G/c^2=1.856\times10^{-27}\,\mathrm{cm\,g^{-1}}$,
$T^{00}=\rho$.
\end{itemize}
Les deux sont \emph{physiquement équivalentes} (le signe et la puissance de $c$ sont absorbés
par la définition de $T_{\mu\nu}$), mais sources fréquentes de malentendus de signe.

\textbf{Schwarzschild extérieure} (vide, symétrie sphérique) :
\begin{equation}
\dd s^2=\Bigl(1-\tfrac{2GM}{c^2r}\Bigr)c^2\dd t^2
-\frac{\dd r^2}{1-\tfrac{2GM}{c^2r}}-r^2\dd\Omega^2,
\qquad r_S=\frac{2GM}{c^2}.
\label{eq:schwext}
\end{equation}
Pour $M<0$, on remplace $1-\tfrac{2GM}{c^2r}$ par $1+\tfrac{2G\abs{M}}{c^2r}$ : \emph{répulsion}
\eqpaper{79}.

\paragraph{Dérivation de Hilbert.} Action d'Einstein–Hilbert
$S_{\mathrm{EH}}=\frac{1}{2\chi}\int R\sqrt{\abs{g}}\dx + S_{\text{mat}}$.
Les trois lemmes variationnels (à maîtriser absolument, ils servent tels quels dans PMZ24
\eqpaper{81}) :

\begin{lemme}[Variations fondamentales]\label{lem:var}
\begin{align}
\delta\sqrt{\abs{g}} &= \half\sqrt{\abs{g}}\,g^{\mu\nu}\delta g_{\mu\nu}
                     = -\half\sqrt{\abs{g}}\,g_{\mu\nu}\delta g^{\mu\nu}, \label{eq:var1}\\
\delta g^{\mu\nu} &= -g^{\mu\alpha}g^{\nu\beta}\delta g_{\alpha\beta}, \label{eq:var2}\\
\delta R &= R_{\mu\nu}\,\delta g^{\mu\nu}+g^{\mu\nu}\delta R_{\mu\nu},
\quad\text{(2$^{\text e}$ terme = divergence totale, Palatini).} \label{eq:var3}
\end{align}
\end{lemme}

\begin{proof}[Esquisse]
\eqref{eq:var1} : $\delta\det g=\det g\,g^{\mu\nu}\delta g_{\mu\nu}$ (formule de Jacobi),
puis $\delta\sqrt{\abs{g}}=\tfrac{1}{2\sqrt{\abs{g}}}\,\delta\abs{g}$.
\eqref{eq:var2} : différentier $g^{\mu\rho}g_{\rho\nu}=\delta^\mu_\nu$.
\eqref{eq:var3} : $g^{\mu\nu}\delta R_{\mu\nu}=\nabla_\sigma w^\sigma$ avec
$w^\sigma=g^{\mu\nu}\delta\Gamma^\sigma_{\mu\nu}-g^{\sigma\nu}\delta\Gamma^\mu_{\mu\nu}$,
qui s'intègre en terme de bord.
\end{proof}

En insérant \eqref{eq:var1}--\eqref{eq:var3} dans $\delta S_{\mathrm{EH}}=0$, et en définissant
\begin{equation}
T_{\mu\nu}=-\frac{2}{\sqrt{\abs{g}}}\frac{\delta(\sqrt{\abs{g}}\,\mathcal L_{\text{mat}})}{\delta g^{\mu\nu}},
\label{eq:Tdef}
\end{equation}
on retrouve \eqref{eq:einstein}.

%==============================================================================
\section{Partie II — Géométrie symplectique de Souriau et groupes dynamiques}
\label{part:souriau}

\subsection{Groupes et algèbres de Lie}

Un \emph{groupe de Lie} $G$ est une variété munie d'une structure de groupe $C^\infty$. Son
\emph{algèbre de Lie} $\mathfrak g=T_eG$ est l'espace tangent à l'identité, munie du crochet
$[\cdot,\cdot]$. Pour Poincaré, $\dim\mathfrak g=10$, élément générique \eqpaper{14} :
\begin{equation}
Z=\begin{pmatrix}\Lambda & \Gamma\\ 0 & 0\end{pmatrix},\quad
\Lambda^{\mathsf T}\eta=-\eta\Lambda,\ \Gamma\in\mathbb R^{1,3}.
\end{equation}

\textbf{Actions adjointe et coadjointe.} $\Ad_g X=gXg^{-1}$ sur $\mathfrak g$ ; sur le dual
$\mathfrak g^\ast$ (l'\emph{espace des moments}), en représentation matricielle \eqpaper{19} :
\begin{equation}
\mu'=g\,\mu\,g^{\mathsf T}.
\label{eq:coadj}
\end{equation}

\subsection{Application moment et espace des moments}

PMZ24 représente $\mu\in\mathfrak g^\ast$ par une matrice antisymétrique $5\times5$
\eqpaper{20--21} :
\begin{equation}
\mu=\begin{pmatrix} M & -P\\ P^{\mathsf T} & 0\end{pmatrix},\quad
M=\begin{pmatrix}
0 & -l_z & l_y & f_x\\
l_z & 0 & -l_x & f_y\\
-l_y & l_x & 0 & f_z\\
-f_x & -f_y & -f_z & 0
\end{pmatrix},\quad
P=\begin{pmatrix}E\\ p_x\\ p_y\\ p_z\end{pmatrix}.
\end{equation}
Les dix coordonnées du \emph{torseur} \eqpaper{26} :
$\mu=\{\,\vb l\ (\text{spin}),\ \vb g\ (\text{barycentre}),\ \vb p\ (\text{impulsion}),\ E\,\}$.

\subsection{Action coadjointe et inversion du temps : le théorème pivot}

Le produit \eqref{eq:coadj} avec $g=\begin{psmallmatrix}L&C\\0&1\end{psmallmatrix}$ donne
\eqpaper{22--25} :
\begin{align}
M'&=LML^{\mathsf T}+CP^{\mathsf T}L^{\mathsf T}-LPC^{\mathsf T}, \label{eq:Mprime}\\
P'&=LP. \label{eq:Pprime}
\end{align}

\begin{theoreme}[T-symétrie $\Rightarrow$ inversion d'énergie]
\label{th:Tinv}
Pour le renversement temporel pur $L=L_t=\diag(-1,1,1,1)$, $C=0$ \eqpaper{28},
\begin{equation}
P'=L_t P=(-E,\,p_x,\,p_y,\,p_z)\qquad\eqpaper{29}.
\end{equation}
Donc $E\to-E$, et par le Théorème~\ref{th:masse}, $m\to-m$ \eqpaper{30}.
\end{theoreme}

C'est l'articulation logique centrale : \emph{le second feuillet de Sakharov est peuplé de
particules à énergie et masse négatives}. (Résultat formel correct ; portée physique
contestée.)

\subsection{Groupe de Kaluza et charge électrique}

On élargit l'espace-temps par une $5^{\text e}$ dimension $\zeta$ compactifiée, de Gram
\eqpaper{33}
\begin{equation}
\Gamma=\begin{pmatrix} G & 0\\ 0 & -1\end{pmatrix},\qquad
\dd s^2=\dd t^2-\dd x^2-\dd y^2-\dd z^2-\dd\zeta^2 \quad\eqpaper{34}.
\end{equation}
Le groupe de Poincaré s'étend \eqpaper{31} par une translation $\varphi$ selon $\zeta$.
Par le théorème de Noether, cette symétrie produit un invariant scalaire $q$ (charge
électrique) ; le torseur devient $\mu=\{\vb l,\vb g,\vb p,E,q\}$ \eqpaper{36}, et l'action
\eqpaper{40--42} donne $q'=q$, plus \eqref{eq:Mprime}--\eqref{eq:Pprime}.

\begin{remarque}[Subtilité relevée par de Saxcé \cite{deSaxce2024}]
Si l'on impose la stricte préservation de la \emph{métrique hyperbolique} 5D, $q$ \emph{n'est
pas} invariante (elle dépend du repère, en contradiction avec l'expérience). L'invariance est
restaurée en compactifiant $\zeta$ avec un rayon $\omega\to0$ (échelle de Planck), soit une
brisure $4{+}1\to4$. La longueur caractéristique \eqpaper{56--57} est
$\omega\approx 3.782\times10^{-32}\,\mathrm{cm}$, et $\omega/2\pi$ redonne la longueur de Planck.
\end{remarque}

\subsection{Groupe Janus restreint et groupe Janus dynamique}

\textbf{Janus restreint} \eqpaper{43} :
$g=\begin{psmallmatrix}\mu&0&\varphi\\0&L&C\\0&0&1\end{psmallmatrix}$, $\mu=\pm1$,
$L=\lambda L_o$. Action $q'=\mu q$ \eqpaper{45} : $\mu=-1$ renverse $\zeta$ et $q$
($C$-symétrie, matière$\leftrightarrow$antimatière de Dirac).

\textbf{Janus dynamique} \eqpaper{58} (construction-pivot) :
\begin{equation}
g=\begin{pmatrix}\lambda\mu & 0 & \varphi\\ 0 & \lambda L_o & C\\ 0 & 0 & 1\end{pmatrix},
\quad \lambda,\mu\in\{-1,+1\},\ L_o\in\Lor_o.
\label{eq:janusdyn}
\end{equation}
Ici $\lambda$ porte la $PT$-symétrie, $\mu$ la $C$-symétrie. D'où les \textbf{quatre types de
particules} (après \eqpaper{58}) :
\begin{center}
\begin{tabular}{cll}
\toprule
$(\lambda,\mu)$ & masse / énergie & nature\\
\midrule
$(+1,+1)$ & positives & matière (notre univers)\\
$(+1,-1)$ & positives & antimatière de Dirac (positrons)\\
$(-1,+1)$ & négatives & $CPT$-image de l'antimatière\\
$(-1,-1)$ & négatives & antimatière de Feynman\\
\bottomrule
\end{tabular}
\end{center}

En utilisant $L_o^{-1}=GL_o^{\mathsf T}G$ \eqpaper{64}, l'action sur $\mathfrak g$ \eqpaper{60,65}
puis sur les moments \eqpaper{72--74} :
\begin{equation}
\sum_i q_i'=\lambda\mu\sum_i q_i,\qquad
M'=LML^{\mathsf T}-LPC^{\mathsf T}+CP^{\mathsf T}L^{\mathsf T},\qquad P'=LP.
\end{equation}
La « violation apparente » de conservation des charges lorsque $\lambda\mu=-1$ est
l'expression mathématique de la \emph{transition entre feuillets}.

%==============================================================================
\section{Partie III — Topologie du modèle Janus}
\label{part:topo}

PMZ24 propose une lecture purement topologique de $P,T$ via la géométrie projective.
En identifiant les points antipodaux de $S^2$ on obtient $\mathbb P^2$ ; cette identification :
(i) transforme le voisinage d'un méridien en couverture à deux feuillets d'une bande de
Möbius à \emph{un} demi-tour ($P$-symétrie, fig.~5 du papier) ; (ii) transforme le voisinage de
l'équateur en bande de Möbius à \emph{trois} demi-tours ($T$-symétrie, fig.~3) ;
(iii) fait coïncider pôles Nord/Sud (Big~Bang / Big~Crunch), singularités obligatoires sur
toute $S^{2n}$. Remplacer la double singularité par un passage tubulaire donne localement une
\emph{bouteille de Klein} (fig.~4). Le cas physique est $S^4\to\mathbb P^4$. Référence
historique : la surface de Boy \cite{Boy1903}, immersion de $\mathbb P^2$ dans $\mathbb R^3$.

\begin{remarque}
Cette partie est \emph{illustrative} : aucune dynamique des feuillets ne s'en déduit
quantitativement. Elle motive l'émergence topologique de $P,T$ mais n'entre pas dans les
équations de champ.
\end{remarque}

%==============================================================================
\section{Partie IV — Les équations de champ bimétriques}
\label{part:champ}

\subsection{Le problème de Bondi (1957) et le runaway effect}

\begin{theoreme}[Runaway effect, Bondi \cite{Bondi1957}]
Dans la RG \emph{mono-métrique}, $+m$ attire tout, $-m$ repousse tout. Une paire $(+m,-m)$
de masses opposées s'auto-accélère indéfiniment dans la même direction \emph{sans apport
d'énergie} (l'énergie cinétique totale $\half(m+(-m))v^2$ reste nulle), violant l'action--réaction.
\end{theoreme}

Dans PMZ24 \S7 \eqpaper{75--79} : une seule métrique $g_{\mu\nu}$ et une source unique
$T^{(+)}_{\mu\nu}+T^{(-)}_{\mu\nu}$ ne fournissent \emph{qu'une} famille de géodésiques, suivie
par toutes les particules-tests. La solution Janus consiste à donner aux masses négatives une
\emph{seconde} famille de géodésiques (celles de $\gbar$).

\subsection{L'action à deux couches et sa variation \eqpaper{80--91}}

\begin{equation}
\boxed{\;
A=\int\!\Bigl(\frac{R}{2\chi}+S+S'\Bigr)\sqrt{\abs{g}}\dx
+\int\!\Bigl(\frac{\kappa\Rbar}{2\chibar}+\Sbar+\Sbar'\Bigr)\sqrt{\abs{\gbar}}\dx
\;}
\label{eq:action}
\end{equation}
où $S,\Sbar$ sont les sources propres aux deux feuillets, $S',\Sbar'$ les termes
d'interaction croisés, et $\kappa=\pm1$. La variation \emph{indépendante} en $\delta g^{\mu\nu}$
et $\delta\gbar^{\mu\nu}$ utilise le Lemme~\ref{lem:var}. On définit \eqpaper{84--87} :
\begin{align}
T_{\mu\nu}&=-\frac{2}{\sqrt{\abs{g}}}\frac{\delta(\sqrt{\abs{g}}\,S)}{\delta g^{\mu\nu}}
=-2\frac{\delta S}{\delta g^{\mu\nu}}+g_{\mu\nu}S,
&\Tbar_{\mu\nu}&=-\frac{2}{\sqrt{\abs{\gbar}}}\frac{\delta(\sqrt{\abs{\gbar}}\,\Sbar)}{\delta\gbar^{\mu\nu}},\\
\Tp_{\mu\nu}&=-\frac{2}{\sqrt{\abs{\gbar}}}\frac{\delta(\sqrt{\abs{g}}\,S')}{\delta g^{\mu\nu}},
&\Tbarp_{\mu\nu}&=-\frac{2}{\sqrt{\abs{g}}}\frac{\delta(\sqrt{\abs{\gbar}}\,\Sbar')}{\delta\gbar^{\mu\nu}}.
\end{align}
Les tenseurs croisés se réécrivent \eqpaper{88--89} :
$\sqrt{\abs{\gbar}/\abs{g}}\;\Tp_{\mu\nu}=-2\,\delta S'/\delta g^{\mu\nu}+g_{\mu\nu}S'$.

\begin{theoreme}[Système couplé, $\kappa=-1$]\label{th:couple}
L'annulation locale \eqpaper{82--83} de la variation, avec le choix $\kappa=-1$, donne
\begin{align}
R_{\mu\nu}-\half R\,g_{\mu\nu}
&=\chi\Bigl[\,T_{\mu\nu}+\sqrt{\tfrac{\abs{\gbar}}{\abs{g}}}\,\Tp_{\mu\nu}\Bigr] &\eqpaper{90},
\label{eq:eq90}\\
\Rbar_{\mu\nu}-\half\Rbar\,\gbar_{\mu\nu}
&=-\chibar\Bigl[\,\Tbar_{\mu\nu}+\sqrt{\tfrac{\abs{g}}{\abs{\gbar}}}\,\Tbarp_{\mu\nu}\Bigr] &\eqpaper{91}.
\label{eq:eq91}
\end{align}
\end{theoreme}

Le signe « $-$ » de \eqref{eq:eq91} (induit par $\kappa=-1$) produira en limite newtonienne la
loi « même signe attire / signes opposés se repoussent ».

\subsection{Solution FLRW homogène isotrope \eqpaper{92--96}}

Deux métriques FLRW de temps cosmique commun $x^0$ \eqpaper{92} :
\begin{equation}
\dd s^2=\dd x^{0\,2}-a^2(x^0)\Bigl[\tfrac{\dd u^2}{1-ku^2}+u^2\dd\Omega^2\Bigr],\quad
\dd\bar s^2=\dd x^{0\,2}-\abar^2(x^0)\Bigl[\tfrac{\dd u^2}{1-\bar k u^2}+u^2\dd\Omega^2\Bigr].
\end{equation}
Déterminants $g=-a^6\sin^2\theta$, $\gbar=-\abar^6\sin^2\theta$ \eqpaper{93}. L'injection dans
\eqref{eq:eq90}--\eqref{eq:eq91} et la compatibilité des deux temps imposent la
\emph{conservation d'énergie généralisée} \eqpaper{94} :
\begin{equation}
\boxed{\,\rho c^2 a^3+\rhobar\cbar^2\abar^3=E=\text{cst}\,}.
\label{eq:conserv}
\end{equation}
Pour la solution exacte de poussière $k=\bar k=-1$ \eqpaper{95}, les équations d'accélération
\eqpaper{96} :
\begin{equation}
a^2\,\ddot a=-\frac{4\pi G}{c^2}\,E,\qquad
\abar^2\,\ddot{\abar}=+\frac{4\pi G}{\cbar^2}\,E.
\label{eq:accel}
\end{equation}

\begin{remarque}[Interprétation]
Si $E<0$ (énergie totale dominée par la masse négative), alors $\ddot a>0$ : \textbf{expansion
accélérée sans constante cosmologique}. La masse négative joue \emph{simultanément} matière
noire (cohésion interne) et énergie noire (répulsion du feuillet~$+$).
\end{remarque}

\textbf{Validation observationnelle.} D'Agostini \& Petit \cite{DAgostini2018} ajustent sur
l'échantillon JLA de 740 SN~Ia \cite{Betoule2014}, obtenant
\begin{equation}
q_0=-0.087\pm0.015,\qquad \chi^2/\mathrm{d.o.f.}=657/738,
\end{equation}
quasi indiscernable de $\Lambda$CDM plat ($\Omega_m=0.295\pm0.034$, $\Omega_\Lambda=0.705$)
sur $0<z<1.4$.

\subsection{Métriques de Schwarzschild positive et négative \eqpaper{109--137}}

\paragraph{Cas 1 : masse positive (\S12.1).}
On pose $\varepsilon R_S=\varepsilon\,2GM/c^2$, $\varepsilon$ petit \eqpaper{109}, métriques
$\dd s^2=e^\nu\dd x^{0\,2}-e^\lambda\dd r^2-r^2\dd\Omega^2$ \eqpaper{110}. Avec
$e^{-\lambda}=1-2m(r)/r$ \eqpaper{116}, le calcul standard \cite[éq.~14.18]{Adler1975} donne
\begin{equation}
m(r)=\frac{G}{c^2}\int_0^r 4\pi r'^2\rho\,\dd r'=\frac{4}{3}\pi r^3\rho\,\frac{G}{c^2}
\quad\eqpaper{117}.
\end{equation}
Avec le tenseur source $T^0_0=\rho c^2$, $T^i_i=-\varepsilon p$ \eqpaper{115}, on obtient la
TOV \eqpaper{118} :
\begin{equation}
\frac{1}{c^2}\dv{p}{r}
=-\frac{\bigl(m+4\pi\varepsilon Gpr^3/c^4\bigr)\bigl(\rho+\varepsilon p/c^2\bigr)}{r(r-2m\varepsilon)}.
\label{eq:tov118}
\end{equation}

\begin{lemme}[Limite newtonienne $\Rightarrow$ Euler]
Quand $\varepsilon\to0$ dans \eqref{eq:tov118} \eqpaper{119} :
\begin{equation}
\dv{p}{r}=-\frac{G\rho\,\mu(r)}{r^2},\qquad \mu(r)=\frac{4}{3}\pi r^3\rho,
\end{equation}
soit l'équilibre hydrostatique newtonien (relation d'Euler).
\end{lemme}

Métrique intérieure \eqpaper{120}, $\hat r=\sqrt{3c^2/(8\pi G\rho)}$ \eqpaper{122}, raccordée à
\eqref{eq:schwext}.

\textbf{Choix du tenseur d'interaction} \eqpaper{123} : $T'^0_0=\rho c^2$, $T'^i_i=+\varepsilon p$
(signe $+$ au lieu de $-$). La TOV correspondante \eqpaper{125} :
\begin{equation}
\frac{1}{c^2}\dv{p}{r}
=-\frac{\bigl(m-4\pi\varepsilon Gpr^3/c^4\bigr)\bigl(\rho-\varepsilon p/c^2\bigr)}{r(r+2m\varepsilon)}.
\end{equation}
\emph{Les deux TOV \eqref{eq:tov118} et \eqpaper{125} ont la même limite d'Euler quand
$\varepsilon\to0$} : c'est l'argument de cohérence asymptotique. La métrique intérieure du
feuillet négatif \eqpaper{126--127} donne une \emph{répulsion} sur une masse négative-test par
une distribution positive.

\paragraph{Cas 2 : masse négative (\S12.2).}
Rôles inversés \eqpaper{128--137} : une masse~$+$ subit une répulsion d'une masse~$-$
\eqpaper{132--133} ; une masse~$-$ subit une attraction d'autres masses~$-$ \eqpaper{136--137}.

\begin{theoreme}[Loi d'interaction synthétique]
\begin{equation*}
(+,+):\text{attraction};\quad(-,-):\text{attraction};\quad(+,-)\text{ ou }(-,+):\text{répulsion}.
\end{equation*}
Le runaway de Bondi est éliminé : les masses opposées se repoussent mutuellement
(action--réaction respectée).
\end{theoreme}

\subsection{Conséquences observationnelles}

\begin{enumerate}[leftmargin=1.6em]
\item \textbf{Structure lacunaire} (fig.~13 ; \cite{Petit1995}) : conglomérats négatifs
confinant la matière positive en « bulles de savon ».
\item \textbf{Dipole repeller} \cite{Hoffman2017} : vide à $\sim$600~M~al, identifié comme
agrégat de masse négative ; bulk flow anti-aligné jusqu'à $16000\pm4500\,\mathrm{km/s}$,
contribution $\approx$ moitié au dipôle CMB $V_{\mathrm{CMB}}=631\pm20\,\mathrm{km/s}$.
\item \textbf{Formation précoce des galaxies} (JWST, \cite{Ferreira2022}) : compression latérale
en plaques minces à refroidissement rapide.
\item \textbf{Prédiction nouvelle} : lentille gravitationnelle \emph{négative} dans les grands
vides (atténuation des magnitudes d'arrière-plan), falsifiable (Euclid, LSST, Roman).
\item \textbf{Tension $H_0$} : valeur « directe » $\sim70$ vs CMB $\Lambda$CDM
$67.4\pm0.5\,\mathrm{km\,s^{-1}\,Mpc^{-1}}$ \cite{Planck2020}.
\end{enumerate}

%==============================================================================
\section{Partie V — Évaluation critique}
\label{part:critique}

\subsection{Cadre général : bigravité mainstream vs Janus}

La gravité bimétrique non-linéaire moderne est dominée par Hassan–Rosen
\cite{HassanRosen2012}, bâtie sur dRGT, d'action
\begin{equation}
S_{\mathrm{HR}}=M_g^2\!\int\!\sqrt{\abs{g}}\,R(g)
+M_f^2\!\int\!\sqrt{\abs f}\,R(f)
+2m^2M_{\text{eff}}^2\!\int\!\sqrt{\abs{g}}\sum_{n}\beta_n\,e_n\!\bigl(\sqrt{g^{-1}f}\,\bigr),
\end{equation}
$e_n$ = polynômes symétriques élémentaires des valeurs propres de $\sqrt{g^{-1}f}$. Cette
construction garantit l'absence du fantôme de Boulware–Deser.

\textbf{Différences avec Janus :} (i) HR a des interactions \emph{algébriques non
dérivatives} entre $g$ et $f$, Janus n'a aucune interaction algébrique (le couplage passe par
$T',\Tbarp$) ; (ii) HR distingue une métrique de référence, Janus fait sourcer les \emph{mêmes}
particules dans les deux équations ; (iii) HR n'a pas de masse négative, Janus en fait son
ingrédient principal. Revue de référence : \cite{SchmidtMay2016}.

\subsection{La critique de Thibault Damour (2019, 2022)}
\label{sec:damour}

\textbf{Références :} \cite{Damour2019,Damour2022}.

\textbf{Argument central.} Les deux Bianchi imposent
$\nabla_\mu E^{(+)\mu\nu}\equiv0$ et $\nablabar_\mu E^{(-)\mu\nu}\equiv0$. Pour un système réduit
où ne subsiste que la matière positive, le système Janus impose à la \emph{même} matière
\begin{equation}
\nabla_\mu T^{\mu\nu}=0 \quad\text{ET}\quad \nablabar_\mu \Tbar^{\mu\nu}=0,
\quad\text{avec } \Tbar_{\mu\nu}\equiv-\tfrac{w}{\bar w}T_{\mu\nu}.
\end{equation}
En approximation newtonienne quasi-statique, $g_{00}=-(1-2U/c^2)$, $\Delta U=-4\pi G\rho$,
et par symétrie formelle $\bar U=-U$. L'équilibre hydrostatique issu des deux conservations
donne deux gradients de pression \emph{incompatibles} :
\begin{equation}
\partial_i p=+\rho\,\partial_i U \ \text{(correct)} \qquad\text{vs}\qquad
\partial_i p=-\rho\,\partial_i U \ \text{(absurde)}.
\end{equation}

\textbf{Conclusion de Damour} (verbatim) : « \emph{Les équations de champ du `modèle Janus'\dots
constituent un système incohérent d'équations.} » La version 2022 montre que l'amélioration
2019 réduit la discrepance ($\sim$200\,\% $\to\sim$20\,\% pour une étoile à neutrons) mais
ne l'élimine pas, « \emph{both at the Newtonian \dots and at the post-Newtonian level} ».

\textbf{Réponse implicite de PMZ24 et évaluation neutre.} PMZ24 introduit les facteurs
$\sqrt{\abs{\gbar}/\abs{g}}$ et des tenseurs $T',\Tbarp$ \emph{distincts} de $T$, et ne réclame
qu'une cohérence \emph{asymptotique} ($\varepsilon\to0$). Mais les Bianchi \eqpaper{97--104}
impliquent toujours $\nabla_\mu T^{\mu\nu}=0$ \emph{et} $\nablabar_\mu\Tbar^{\mu\nu}=0$ ; les
auteurs admettent eux-mêmes (p.~17) que la cohérence n'est garantie qu'en limite newtonienne
via $\sqrt{\abs{g}/\abs{\gbar}}\to1$. La stratégie 2024 \emph{atténue} mais ne \emph{résout pas}
l'objection de fond.

\subsection{Autres critiques publiques}

\begin{itemize}[leftmargin=1.4em]
\item \textbf{A.~Riazuelo (IAP, 2006)} : critique de Petit~1995 (signe de courbure, définition de
$T_{\mu\nu}$, problème de l'horizon) ; ne porte pas sur la version 2024.
\item \textbf{L.~J.~Bourhis (2022)} \cite{Bourhis2022} : restreint l'argument de Damour au cas
mono-sectoriel ; diffusion limitée, non revue par les pairs.
\item \textbf{Boulanger–Damour–Gualtieri–Henneaux (2001)} : théorème no-go général pour
théories multi-spin-2 en interaction ; Janus n'entre dans aucune classe ghost-free.
\item \textbf{Hohmann \& Wohlfarth (2009)} \cite{Hohmann2009} : no-go pour bigravité à masses
$\pm$ ; Hossenfelder \cite{Hossenfelder2009} en conteste l'applicabilité à son modèle, pas à Janus.
\item \textbf{Socas-Navarro (2019)} \cite{SocasNavarro2019} et \textbf{Stepanian} \cite{Stepanian2019}
critiquent le modèle Farnes à masse négative ; arguments en partie transposables, Janus
échappant au runaway \emph{à courte distance} par construction.
\end{itemize}

\subsection{Controverse éditoriale (EPJC)}

À mai~2026, \emph{aucune} action éditoriale (rétractation, expression of concern, note) n'a été
émise par EPJC sur PMZ24 : la page Springer affiche « Open access, Published 26~November 2024 »
sans bannière post-publication ; aucune entrée Retraction~Watch. Par contraste, EPJC a émis une
Expression of Concern puis une Retraction Note sur un autre article (Momeni 2025, pour
fabrication de références par IA), preuve que le journal mobilise ces outils quand justifié.
L'absence d'action sur PMZ24 est donc un résultat \emph{informatif} (et non une simple omission).

\textbf{Statut sociologique.} Petit (né 1937, ancien directeur de recherche CNRS) est une figure
hétérodoxe ; le conflit ancien avec Damour, Riazuelo et l'IAP est documenté publiquement. Le
débat est en partie devenu \emph{méta-scientifique}.

\subsection{Le modèle est-il falsifiable ?}

\textbf{Post-dictions} : expansion accélérée, grands vides, formation rapide des galaxies,
dipole repeller. \textbf{Prédictions véritablement nouvelles} : lentille négative dans les
vides (testable Euclid/LSST/Roman) ; absence de structure dans le feuillet négatif (non
directement observable) ; $q_0=-0.087$ (testable avec un plus grand échantillon SN). La
falsifiabilité poppérienne réelle repose surtout sur la lentille négative, dont la précision
quantitative n'est pas pleinement spécifiée.

\subsection{Synthèse}

\textbf{Cohérence interne} : la construction de groupes (Poincaré, Kaluza, Janus) est
formellement rigoureuse ; l'identification $T$-symétrie $\leftrightarrow$ inversion de masse est
correcte ; la loi « même signe attire » répond constructivement à Bondi.

\textbf{Contestation externe} : (1) l'objection des deux Bianchi n'est levée qu'en limite
newtonienne ; (2) les Lagrangiens d'interaction $S',\Sbar'$ ne sont jamais explicités (théorie
définie par ses variations) ; (3) le modèle n'est dans aucune classe ghost-free ; (4) la règle
« chaque masse suit sa propre métrique » n'est pas dérivée d'un principe variationnel cohérent ;
(5) la cosmologie quantitative (Planck/DES/Euclid) ne cite quasiment jamais le modèle.

%==============================================================================
\section{Programme de travail recommandé}

\begin{description}[leftmargin=0pt,style=nextline]
\item[Étape 1 — Prérequis (2–4 sem.)] Carroll \cite{Carroll2004} ch.~1–6 ; Adler–Bazin–Schiffer
\cite{Adler1975} ch.~5–7. Maîtriser Levi-Civita, dérivée covariante, Riemann/Ricci, Bianchi,
action d'Einstein–Hilbert.
\item[Étape 2 — Souriau (2 sem.)] \cite{Souriau1970} ch.~13–14 ; \cite{deSaxceMarle2023}.
Refaire \eqpaper{22--25} et \eqpaper{65}.
\item[Étape 3 — Calculs du papier (1 mois)] Reproduire l'action coadjointe \eqpaper{22--26},
les groupes \eqpaper{39--74}, la variation $\delta A=0$ \eqpaper{80--91}, la solution FLRW
\eqpaper{92--96}, les TOV \eqpaper{118} et \eqpaper{125}.
\item[Étape 4 — Étude critique (2 sem.)] Damour \cite{Damour2019,Damour2022} intégralement ;
Hassan–Rosen \cite{HassanRosen2012} ; revue \cite{SchmidtMay2016} ; Bondi \cite{Bondi1957}.
\end{description}

\textbf{Jalons de révision de position.} (i) Euclid (DR1 prévu 21~oct.~2026) confirmant une
lentille négative $>3\sigma$ $\Rightarrow$ crédibilité accrue. (ii) Une dérivation lagrangienne
revue par les pairs (hors PMZ) répondant aux deux Bianchi $\Rightarrow$ incohérence levée.
(iii) $q_0$ contraint $<0.01$ et $\neq-0.087$ $\Rightarrow$ falsification. (iv) Expression of
concern/rétractation EPJC $\Rightarrow$ poids éditorial affaibli.

%==============================================================================
\section{Réserves méthodologiques (\emph{caveats})}

\begin{enumerate}[leftmargin=1.6em,nosep]
\item La variation $\delta A=0$ \eqpaper{81--91} n'est pas détaillée pas-à-pas dans le papier ;
à refaire soi-même via le Lemme~\ref{lem:var}.
\item La forme explicite de $S',\Sbar'$ n'est jamais donnée : $T',\Tbarp$ sont définis par leurs
variations (boîte noire prescriptive).
\item La validité asymptotique des Bianchi \eqpaper{99--104} ne garantit pas la cohérence
globale ; aucun théorème d'existence/unicité en régime fort.
\item L'invariance électrique \eqpaper{56} repose sur une formule heuristique
(Souriau~1964 p.~412), rigoureusement justifiée seulement via de Saxcé~2024.
\item Le statut topologique ($S^4\to\mathbb P^4$, surface de Boy) est illustratif.
\item Aucune publication post-2024 revue par les pairs ne discute substantiellement PMZ24
(hors extensions des auteurs eux-mêmes).
\item La convention $\chi=-8\pi G/c^2$ \eqpaper{57} diffère de celle de Damour ($+8\pi G/c^4$) :
équivalentes mais sources de confusions de signe.
\item Les « prédictions » sont majoritairement des post-dictions.
\item Sociologie de la controverse : ne s'appuyer que sur sources primaires/critiques signées.
\item $q_0=-0.087\pm0.015$ provient du seul ajustement D'Agostini \& Petit~2018 sur la donnée
JLA de Betoule et al.~2014 (même donnée, deux paramétrisations).
\end{enumerate}

%==============================================================================
\begin{thebibliography}{99}
\addcontentsline{toc}{section}{Références}

\bibitem{PMZ2024} J.-P.~Petit, F.~Margnat, H.~Zejli,
\emph{A bimetric cosmological model based on Andreï Sakharov's twin universe approach},
Eur. Phys. J. C \textbf{84}:1226 (2024).
DOI~\href{https://doi.org/10.1140/epjc/s10052-024-13569-w}{10.1140/epjc/s10052-024-13569-w}.

\bibitem{Souriau1970} J.-M.~Souriau,
\emph{Structure des systèmes dynamiques}, Dunod, Paris (1970).
Trad.~angl.~\emph{Structure of Dynamical Systems}, Birkhäuser (1997).

\bibitem{Souriau1964} J.-M.~Souriau,
\emph{Géométrie et Relativité}, Hermann, Paris (1964).

\bibitem{deSaxceMarle2023} G.~de Saxcé, C.-M.~Marle,
\emph{Presentation of Jean-Marie Souriau's book ``Structure des systèmes dynamiques''} (2023),
\href{https://arxiv.org/abs/2306.03106}{arXiv:2306.03106}.

\bibitem{deSaxce2024} G.~de Saxcé,
\emph{Which symmetry group for elementary particles with an electric charge today and in the past?} (2024),
\href{https://arxiv.org/abs/2403.14846}{arXiv:2403.14846}; publié J. Geom. Phys. \textbf{213}, 105491 (2025).

\bibitem{Bondi1957} H.~Bondi,
\emph{Negative Mass in General Relativity}, Rev. Mod. Phys. \textbf{29}(3), 423--428 (1957).
DOI~\href{https://doi.org/10.1103/RevModPhys.29.423}{10.1103/RevModPhys.29.423}.

\bibitem{Sakharov1967} A.~D.~Sakharov,
\emph{Violation of CP invariance, C asymmetry, and baryon asymmetry of the universe},
JETP Lett. \textbf{5}, 24 (1967).

\bibitem{DAgostini2018} G.~D'Agostini, J.-P.~Petit,
\emph{Constraints on Janus cosmological model from recent observations of supernovae type Ia},
Astrophys. Space Sci. \textbf{363}:139 (2018).
DOI~\href{https://doi.org/10.1007/s10509-018-3365-3}{10.1007/s10509-018-3365-3}.

\bibitem{Betoule2014} M.~Betoule et al.,
\emph{Improved cosmological constraints from a joint analysis of the SDSS-II and SNLS supernova samples},
Astron. Astrophys. \textbf{568}, A22 (2014).
DOI~\href{https://doi.org/10.1051/0004-6361/201423413}{10.1051/0004-6361/201423413}.

\bibitem{Petit1995} J.-P.~Petit,
\emph{Twin Universes Cosmology}, Astrophys. Space Sci. \textbf{226}, 273--307 (1995).
DOI~\href{https://doi.org/10.1007/BF00627375}{10.1007/BF00627375}.

\bibitem{Hoffman2017} Y.~Hoffman, D.~Pomarède, R.~B.~Tully, H.~M.~Courtois,
\emph{The dipole repeller}, Nat. Astron. \textbf{1}, 0036 (2017).
DOI~\href{https://doi.org/10.1038/s41550-016-0036}{10.1038/s41550-016-0036}.

\bibitem{Ferreira2022} L.~Ferreira et al.,
\emph{Panic! At the Disks: \dots galaxy structure at $z>3$ with JWST in the SMACS 0723 field},
Astrophys. J. Lett. \textbf{938}(1), L2 (2022).
DOI~\href{https://doi.org/10.3847/2041-8213/ac947c}{10.3847/2041-8213/ac947c}.

\bibitem{Planck2020} Planck Collaboration,
\emph{Planck 2018 results. VI. Cosmological parameters}, Astron. Astrophys. \textbf{641}, A6 (2020),
\href{https://arxiv.org/abs/1807.06209}{arXiv:1807.06209}.

\bibitem{Boy1903} W.~Boy,
\emph{Über die Curvatura integra und die Topologie geschlossener Flächen},
Math. Ann. \textbf{57}(2), 151--184 (1903).
DOI~\href{https://doi.org/10.1007/BF01444342}{10.1007/BF01444342}.

\bibitem{Adler1975} R.~Adler, M.~Bazin, M.~Schiffer,
\emph{Introduction to General Relativity}, McGraw-Hill, New York (1975).

\bibitem{Carroll2004} S.~Carroll,
\emph{Spacetime and Geometry}, Pearson (2004) ;
notes \href{https://arxiv.org/abs/gr-qc/9712019}{arXiv:gr-qc/9712019}.

\bibitem{HassanRosen2012} S.~F.~Hassan, R.~A.~Rosen,
\emph{Bimetric Gravity from Ghost-free Massive Gravity}, JHEP \textbf{02}, 126 (2012),
\href{https://arxiv.org/abs/1109.3515}{arXiv:1109.3515}.

\bibitem{SchmidtMay2016} A.~Schmidt-May, M.~von Strauss,
\emph{Recent developments in bimetric theory}, J. Phys. A \textbf{49}, 183001 (2016),
\href{https://arxiv.org/abs/1512.00021}{arXiv:1512.00021}.

\bibitem{Damour2019} T.~Damour,
\emph{Sur le ``modèle Janus'' de J.~P.~Petit}, IHES, 4~janvier 2019.
\url{https://www.ihes.fr/~damour/publications/JanusJanvier2019-1.pdf}.

\bibitem{Damour2022} T.~Damour,
\emph{Incohérence Physique et Mathématique du ``modèle Janus-2019'' de J.~P.~Petit et coll.},
IHES, 28~décembre 2022.
\url{https://www.ihes.fr/~damour/publications/JanusDecembre2022c.pdf}.

\bibitem{Bourhis2022} L.~J.~Bourhis,
\emph{Le modèle Janus de J.P.~Petit est irrémédiablement faux}, GitHub Gist (2022).
\url{https://gist.github.com/luc-j-bourhis/c6d395b966e47031654426292e8d898e}.

\bibitem{Hohmann2009} M.~Hohmann, M.~N.~R.~Wohlfarth,
\emph{No-go theorem for bimetric gravity with positive and negative mass},
Phys. Rev. D \textbf{80}, 104011 (2009), \href{https://arxiv.org/abs/0908.3384}{arXiv:0908.3384}.

\bibitem{Hossenfelder2009} S.~Hossenfelder,
\emph{Comment on ``No-go theorem for bimetric gravity\dots''} (2009),
\href{https://arxiv.org/abs/0909.2094}{arXiv:0909.2094}.

\bibitem{SocasNavarro2019} H.~Socas-Navarro,
\emph{Can a negative-mass cosmology explain dark matter and dark energy?},
Astron. Astrophys. \textbf{626}, A5 (2019), \href{https://arxiv.org/abs/1902.08287}{arXiv:1902.08287}.

\bibitem{Stepanian2019} A.~Stepanian,
\emph{On the invalidity of ``negative mass'' description of the dark sector},
Mod. Phys. Lett. A (2019), \href{https://arxiv.org/abs/1903.06037}{arXiv:1903.06037}.

\end{thebibliography}

\end{document}